\section{Semantic Concepts}

Now that we have a handle on the truth-functional semantics of TFL, we are
better positioned to put meat on the bones of some informal concepts introduced
a few weeks ago: necessary truth/falsehood, contingent truth/falsehood,
equivalence and joint possibility.

We can use truth tables to define a technical notions that approximate the
concepts of necessary truth and necessary falsity.

\begin{definition}[Tautology]
  \label{def:tautology}
  \ephi is a \emph{tautology} in TFL iff it is true on every valuation.
\end{definition}

\begin{definition}[Contradiction]
  \label{def:contradiction}
  \ephi is a \emph{contradiction} in TFL iff it is false on every valuation.
\end{definition}

For example, $A \eand \enot A$ is a contradiction whereas $A \eor \enot A$ is
a tautology.\footnote{I've only filled in the values for the main operator, but
the remaining columns for these one atom sentences are trivial. Later we will
learn how to omit as much material as possible to save time as space.}

\begin{center}
  \begin{tabular}{c|c|c}
    $A$ & $A \eand \enot A$ & $A \eor \enot A$ \\
    \hline
    T & \bf{F} & \bf{T} \\
    F & \bf{F} & \bf{T}    
  \end{tabular}
\end{center}

Similarly, we can test whether two sentences are equivalent by checking to
see if they have the same truth tables.

\begin{definition}[Equivalence]
  \label{def:equivalence}
  \ephi and \epsi are \emph{equivalent} in TFL iff \ephi and \epsi have the
  same truth value as each other in every valuation.
\end{definition}

So we can easily verify that $\enot(A\eand B)$ is equivalent to $\enot A \eor
\enot B$ since they have the same truth value in every row.

\begin{center}
  \begin{tabular}{c c|c c c c|c c c c c}
    $A$&$B$ & \enot&$(A$&\eand&$B)$ & \enot&$A$&\eor&\enot&$B$ \\
    \hline
    T & T  &  \bf{F} & T & T & T  &  F & T & \bf{F} & F & T \\
    T & F  &  \bf{T} & T & F & F  &  F & T & \bf{T} & T & F \\
    F & T  &  \bf{T} & F & F & T  &  T & F & \bf{T} & F & T \\
    F & F  &  \bf{T} & F & F & F  &  T & F & \bf{T} & T & F \\
    \hline
      &    &  \bf{4} & 1 & 3 & 2  &  3 & 1 & \bf{5} & 4 & 2
  \end{tabular}
\end{center}

And we can just as easily use truth table methods to evaluate the logical
properties of sets of sentences, including arguments.

\begin{definition}[Joint Satisfiability]
  \label{def:jointly-satisfiable}
  The sentences $\ephi_1\dots\ephi_n$ are \emph{jointly satisfiable} in TFL iff
  some valuation makes them all true. Otherwise, they are \emph{jointly
  unsatisfiable}.
\end{definition}

\begin{definition}[Entailment]
  \label{def:entailment}
  The sentences $\ephi_1\dots\ephi_n$ \emph{entail} \epsi in TFL
  ($\ephi_1\dots\ephi_n\eentails\epsi$) iff no valuation makes
  $\ephi_1\dots\ephi_n$ all true but \epsi false.
\end{definition}

TODO: Discussion of entailment and material implication.
\cite[][pp. 88-90]{Mag25}

As we've already seen, the symbolization of English sentences into TFL leaves
much of the linguistic nuance out of the resulting TFL sentence. We can
summarize some of the important shortcomings of TFL symbolization as follows:

\begin{enumerate}
\item TFL does not support the types of inferences depicted in Aristotles
  syllogisms (quantification).
\item TFL insists on `excluding the middle' -- every TFL is either true
  or false, but some adjectives do not quite work like that.
\item TFL comes with a number of `paradoxes' that make the material conditional
  unsatisfying for symbolizing English sentences: one the one hand, negating the
  whole conditional often gives the wrong conclusions about its antecedent
  sentence; on the other hand, negating the consequent sentence opens up the
  possibility of contradiction when the antecedant sentence is false.
\end{enumerate}

TODO: Examples?

\subsection{Truth table shortcuts} \cite[][Chapter 14, pp. 97-101]{Mag25}

The truth table method is tedious to say the least. Even when the number of
atomic sentences makes the overall number of rows manageable, a logician would
need to write dozens of `T' and `F' symbols to completely fill in a table. We
can introduce a number of shortcuts to make the process much more ergunomic.

\begin{enumerate}
\item You do not need to repeat truth values for the atomic letters. They are
  easy enough to read right off the left column.
\item Since \ephi\eor\epsi is true just in case either \ephi or \epsi is true,
  we can skip one disjunct if the other is already marked `T'.
\item Since \ephi\eand\epsi is false just in case either \ephi or \epsi is
  false, we can skip one conjunct if the other is already marked `F'.
\item Since \ephi\eif\epsi is true in case either \epsi is true or \ephi is
  false, we can skip one component if the other is already so marked.
\end{enumerate}

With these simiplified methods, we can much more efficiently evaluate the
validity of TFL arguments. Like the validity of individual sentences, testing
the validity of arguments depends on identifying the load-bearing sentences
in the structure of the argument. Tersely, we look for `bad' lines.

\begin{enumerate}
\item When the conclusion is true, the line is not bad.
\item When any of the premises is false, the line is not bad.
\item Truth tables with no bad lines depict valid arguments.
\end{enumerate}

TODO: Examples??

